\section{Research Questions}
\label{sec:research-questions}

% Our work addresses the overarching problem of how energy efficiency can be made a design time concern for software developers.
To guide our investigation, we formulate the following three research questions. These questions are answered through the analyses presented in Section~\ref{sec:results}.

\textbf{RQ1: How can energy consumption of small code blocks be measured accurately and reproducibly?}  
We evaluate the proposed methodology's measurement granularity, stability and effectiveness for small code blocks.

\textbf{RQ2: What statistical relationships exist between static code features and block-level energy consumption?}  
We investigate whether code structure and metrics are statistically related to energy consumption. This leads to two sub-questions:
\begin{itemize}
  \item \textbf{RQ2a}: Which static code features exhibit associations with block's energy consumption?  
  \item \textbf{RQ2b}: What is the nature of these associations between energy consumption and code features?
  %   \item \textbf{RQ2a}: What kind of code aspects have high impact on block's energy consumption?  
  % \item \textbf{RQ2b}: How do structural and contextual features derived from Abstract Syntax Trees correlate to energy?
\end{itemize}
% These are answered through correlation analyses (Pearson, Spearman, Kendall) and feature importance rankings across features.

\textbf{RQ3: How reliably can code features predict block-level energy consumption?}  
Beyond correlation, we ask whether predictive models trained on static features can predict energy consumption and generalize to unseen code. This leads to two sub-questions:
\begin{itemize}
  \item \textbf{RQ3a:} How accurately can regression models predict absolute block-level energy values?  
  \item \textbf{RQ3b:} How effectively can classification models assign blocks to energy tiers (low, medium, high)?  
\end{itemize}
% These are answered through regression and classification performance metrics, cross-validation stability, and analysis of generalization gaps between training and test sets.

% These questions are answered through the implementation of our Visual Studio Code extension, which highlights energy-intensive blocks inline with source code, enabling targeted refactoring at design time.
